\begin{center}
\textbf{\large{Abstract}}
\end{center}
\addcontentsline{toc}{chapter}{Abstract}

\begin{small}

Database management systems (DBMS) are critical infrastructure components, and their vulnerabilities represent serious security threats. Automated vulnerability detection through fuzzing is a promising approach, but existing grammar-based fuzzers apply uniform random sampling over production rules, wasting energy (time) on low-value syntactic paths and failing to prioritize the structural patterns most likely to expose bugs.

This thesis presents DBMS-Nautilus, a grammar-based greybox fuzzing system targeting SQLite. What distinguishes DBMS-Nautilus from existing fuzzers is its grammar: a context-free grammar of 526 production rules that encodes vulnerability-triggering SQL patterns --- integer overflow in expression evaluation, use-after-free in query planner optimizations, and null pointer dereference in column resolution --- without embedding any proof-of-concept exploit input. The grammar is the primary contribution: it determines what the fuzzer can find, while the fuzzing engine merely navigates within the space the grammar defines.

Experiments across 4 SQLite versions show that DBMS-Nautilus with seed rules successfully rediscovers all 4 target CVEs through compositional grammar generation. Without seed rules, it discovers 13 distinct bug classes across 48 unique stack hashes --- including 4 independent CVE rediscoveries, 6 non-CVE bugs that crash production builds across up to 10 consecutive releases, and 3 sanitizer-only detections. A comparative evaluation against a generic EBNF baseline grammar demonstrates that the proposed grammar discovers $108\times$ more unique root causes per campaign despite $2.1\times$ lower throughput, indicating that grammar design has a larger effect on unique crash diversity than raw execution speed.

\vspace*{1cm}
\textbf{Keywords}: Greybox Fuzzing, DBMS, Grammar-based, SQLite
\end{small}
